El diseño de una plataforma interactiva de programación requiere integrar un entorno capaz de compilar y ejecutar el código fuente de los estudiantes para ofrecer una experiencia pedagógica interactiva basada en la retroalimentación inmediata. Esta necesidad funcional implica habilitar la ejecución de código arbitrario como un servicio el cual expone inevitablemente al servidor anfitrión a múltiples vectores de amenaza como la escalada de privilegios para acceder al sistema de archivos, la exfiltración de credenciales sensibles mediante conexiones de red salientes y los ataques de denegación de servicio (DoS) provocados por bucles infinitos o la creación exponencial de procesos mediante \textit{fork bombs} destinados a agotar los recursos del sistema operativo \cite{sun2018security}. Para mitigar estos riesgos, es necesario implementar un confinamiento estricto de procesos mediante técnicas de \textit{sandboxing} que garanticen que el código no confiable funcione en un entorno estanco y sin privilegios de administración \cite{wan2019practical}.

Frente a estos riesgos las plataformas web modernas descartan las máquinas virtuales tradicionales debido a su alta latencia y exploran alternativas como la ejecución en el cliente mediante \textbf{Pyodide} la cual utiliza \textit{WebAssembly} para procesar el código directamente en el navegador del estudiante. Aunque este enfoque garantiza una seguridad absoluta para el servidor al delegar el cómputo íntegramente al hardware del usuario, presenta desventajas como la elevada carga inicial de librerías pesadas y la dependencia de los recursos locales del dispositivo del alumno lo que dificulta una experiencia de tutoría consistente. Por otro lado están los motores de ejecución efímeros basados en la virtualización ligera del núcleo de Linux donde destaca \textbf{Piston} al orquestar trabajos aislados (\textit{jobs}) mediante el uso coordinado de \textit{namespaces} y \textit{cgroups} para seccionar la visibilidad de red y limitar el consumo de recursos físicos de forma simultánea \cite{rosen2013resource}. Este software protege la infraestructura al ejecutar cada envío en directorios temporales montados sobre la memoria RAM y bajo identidades de usuario huérfanas sin privilegios administrativos garantizando que cualquier proceso sospechoso sea interceptado por un supervisor temporal que termina la ejecución automáticamente antes de comprometer la integridad del sistema global \cite{engineerman_piston}. No obstante es fundamental reconocer que la ejecución de código arbitrario siempre conlleva un peligro por lo que estas capas de blindaje se plantean como una estrategia de defensa en profundidad orientada a mitigar los riesgos conocidos sin ignorar la posibilidad teórica de vulnerabilidades imprevistas en el aislamiento del núcleo.